\documentclass{beamer}
\usepackage{amsmath,amsfonts}
\usetheme{Madrid}

\title{Introduction to Differential Equations}
\author{}
\date{}

\begin{document}

\frame{\titlepage}

\begin{frame}{What is a Differential Equation?}
A \textbf{differential equation} is an equation that relates an unknown function to its derivatives. Goal: find a function that satisfies the equation.

Examples:

% 1. Exponential Growth/Decay
\[
\frac{dy}{dx} = y
\]
The rate of change is equal to the current amount. (e.g., population growth, continuous interest)

% 2. Simple Harmonic Motion
\[
\frac{d^2x}{dt^2} + \omega^2 x = 0
\]
Acceleration is opposite and proportional to displacement. (e.g., a swinging pendulum or a vibrating spring)

% 3. The Heat Equation (Diffusion)
\[
\frac{\partial u}{\partial t} = k \frac{\partial^2 u}{\partial x^2}
\]
How temperature or particle concentration smooths out over time. (e.g., heat spreading through a metal rod)
\end{frame}


\begin{frame}{Newton's Law of Cooling}
How fast does a hot cup of coffee cool down? The rate of cooling is proportional to the temperature difference between the object and its environment.

\vspace{0.5cm}
\textbf{The Equation:}
\[
\frac{dT}{dt} = -k(T - T_{\text{env}})
\]

\begin{itemize}
    \item $T$: Temperature of the object
    \item $T_{\text{env}}$: Ambient temperature (constant)
    \item $k$: Cooling constant (depends on material and surface area)
\end{itemize}

\vspace{0.5cm}
\textbf{How it was derived/measured:} \\
Empirically observed by Isaac Newton in 1701 using a red-hot block of iron. In a modern lab, if you measure $T$ over time and plot $\ln(T - T_{\text{env}})$ against $t$, you get a perfect straight line where the slope is $-k$.
\end{frame}


\begin{frame}{The Simple Pendulum}
What dictates the swing of a grandfather clock? The restoring force pulling the pendulum back to the center depends on gravity and the angle of the swing.

\vspace{0.5cm}
\textbf{The Equation:}
\[
\frac{d^2\theta}{dt^2} + \frac{g}{L}\sin(\theta) = 0
\]

\begin{itemize}
    \item $\theta$: Angle of displacement
    \item $g$: Acceleration due to gravity
    \item $L$: Length of the pendulum
\end{itemize}

\vspace{0.5cm}
\textbf{How it was derived/measured:} \\
Derived directly from Newton's Second Law for rotation ($\tau = I\alpha$). Because $\sin(\theta)$ makes the math non-linear, physicists often measure very small swings where $\sin(\theta) \approx \theta$, turning it into a simple harmonic oscillator.
\end{frame}

\begin{frame}{The RC Circuit \& The Biological Neuron}
How does a cell membrane charge up before firing an action potential? It acts exactly like a resistor and capacitor in parallel.

\vspace{0.5cm}
\textbf{The Equation:}
\[
C_m \frac{dV}{dt} + \frac{V - V_{\text{rest}}}{R_m} = I_{\text{ext}}(t)
\]

\begin{itemize}
    \item $V$: Membrane voltage
    \item $C_m$: Membrane capacitance (lipid bilayer)
    \item $R_m$: Membrane resistance (ion channels)
    \item $I_{\text{ext}}$: External stimulus current
\end{itemize}

\vspace{0.5cm}
\textbf{How it was derived/measured:} \\
Derived theoretically from Kirchhoff's Current Law (conservation of charge). It was famously measured and proven in the 1950s by Hodgkin and Huxley by inserting microscopic electrodes into the giant axons of squids.
\end{frame}

\begin{frame}{Ordinary vs Partial Differential Equations}
\textbf{Ordinary Differential Equations (ODE)}

Unknown function depends on one variable.

Example:

\[
\frac{dy}{dx} = x^2
\]

\textbf{Partial Differential Equations (PDE)}

Unknown function depends on several variables.

Example:

\[
\frac{\partial u}{\partial t} = k \frac{\partial^2 u}{\partial x^2}
\]
\end{frame}

\begin{frame}{Order of a Differential Equation}
The \textbf{order} of a differential equation is the highest derivative appearing in it.

First order:
\[
\frac{dy}{dx} = x
\]

Second order:
\[
\frac{d^2y}{dx^2} + y = 0
\]

Third order:
\[
\frac{d^3y}{dx^3} + y = 0
\]
\end{frame}

\begin{frame}{Linear vs Nonlinear Equations}
\textbf{Linear differential equation}

\[
a_n(x)y^{(n)} + a_{n-1}(x)y^{(n-1)} + \dots + a_1(x)y' + a_0(x)y = f(x)
\]

Example:

\[
y'' + 3y' + 2y = 0
\]

\textbf{Nonlinear examples}

\[
y' = y^2
\]

\[
y'' + y^3 = 0
\]
\end{frame}

\begin{frame}{Homogeneous vs Nonhomogeneous}
\textbf{Linear Homogeneous}

\[
y'' + y = 0
\]

\textbf{Linear Nonhomogeneous}

\[
y'' + y = \sin x
\]

Interpretation:

Homogeneous -- free system

Nonhomogeneous -- driven system
\end{frame}

\begin{frame}{Autonomous vs Nonautonomous}
\textbf{Autonomous (no $t$)}

\[
\frac{dx}{dt} = f(x)
\]

Example:

\[
\frac{dx}{dt} = x(1-x)
\]

\textbf{Non-autonomous}

\[
\frac{dx}{dt} = f(x,t)
\]

Example:

\[
\frac{dx}{dt} = x + \sin t
\]
\end{frame}

\begin{frame}{Systems of Differential Equations}
Example system:

\[
\begin{cases}
\dot{x} = y \\
\dot{y} = -x
\end{cases}
\]

Vector form:

\[
\frac{d}{dt}
\begin{pmatrix}
 x \\
 y
\end{pmatrix}
=
\begin{pmatrix}
 y \\
 -x
\end{pmatrix}
\]
\end{frame}


\begin{frame}{Initial vs Boundary Value Problems}
\textbf{Initial Value Problem (IVP)}
\begin{itemize}
    \item \textbf{The Math:} $y' = y, \quad y(0) = 1$
    \item \textbf{The Physics:} Evolution over time. You know the starting state at $t=0$, and the equation drives it forward. (e.g., radioactive decay, or launching a rocket).
    \item \textbf{Solvability:} \textit{Unambiguous.} As long as the system is well-behaved, theorems guarantee exactly one unique path forward into the future.
\end{itemize}

\vspace{0.4cm}

\textbf{Boundary Value Problem (BVP)}
\begin{itemize}
    \item \textbf{The Math:} $y'' + y = 0, \quad y(0)=0, \quad y(\pi)=0$
    \item \textbf{The Physics:} Spatial constraints. You know the conditions at the physical edges of a domain. (e.g., a guitar string pinned at both ends, or the temperature at the two ends of a metal rod).
    \item \textbf{Solvability:} \textit{Tricky!} BVPs can have one unique solution, no solution at all, or infinitely many. (The math example above has infinitely many solutions: $y = c \sin(x)$).
\end{itemize}
\end{frame}

\begin{frame}{Existence \& Uniqueness: The Picard Operator}
Consider the Initial Value Problem (IVP):
\[
\frac{dy}{dt} = f(t, y) \quad \text{with} \quad y(t_0) = y_0
\]

\textbf{Step 1: Integrate both sides.} \\
If we integrate from $t_0$ to $t$, the Fundamental Theorem of Calculus gives us:
\[
y(t) = y_0 + \int_{t_0}^t f(s, y(s)) \, ds
\]
\textit{Notice: $y$ is on both sides. This is an integral equation!}

\textbf{Step 2: Define the Picard Operator ($T$).} \\
Let $T$ be a "machine" that takes a function $\phi(t)$ and spits out a new function:
\[
(T\phi)(t) = y_0 + \int_{t_0}^t f(s, \phi(s)) \, ds
\]
\textbf{The Goal:} We are looking for a function $y$ such that $Ty = y$. This is a \textbf{fixed point}!
\end{frame}

\begin{frame}{Sketch of Proof: Banach Fixed-Point Theorem}
How do we find the fixed point $Ty = y$? We guess and iterate!

\textbf{The Picard Iteration:}
\begin{enumerate}
    \item Start with a naive guess: $\phi_0(t) = y_0$ (a flat line).
    \item Plug it into the operator: $\phi_1 = T\phi_0$.
    \item Repeat indefinitely: $\phi_{n+1} = T\phi_n$.
\end{enumerate}

\textbf{Why does this work?} \\
If $f(t, y)$ is "well-behaved" (Lipschitz continuous), then $T$ is a \textbf{contraction mapping}. This means every time you apply $T$, the distance between any two functions shrinks.

By the \textbf{Banach Fixed-Point Theorem}, if you keep applying a contraction mapping in a complete space, you are guaranteed to get sucked into one, single, unique fixed point.

\[
\lim_{n \to \infty} \phi_n(t) = y(t) \quad \text{(The unique solution!)}
\]
\end{frame}


\begin{frame}{Exact vs Numerical Solutions}
\textbf{Analytical solution}

Closed form expression.

Example:

\[
y = e^x
\]

\textbf{Numerical methods}

\begin{itemize}
\item Euler method
\item Runge--Kutta
\item Finite difference methods
\item Finite element methods
\end{itemize}
\end{frame}

\begin{frame}{Solving the IVP: Unambiguous Future}
\textbf{Problem:} $y' = y$, with the initial state $y(0) = 1$.

\textbf{Step 1: Find the General Solution} \\
This is a separable equation. We gather variables:
$$\frac{dy}{y} = dx$$

Integrate both sides:
$$\ln|y| = x + C \implies y(x) = Ce^x$$
\textit{(This represents a family of infinite parallel curves.)}

\textbf{Step 2: Apply the Initial Condition} \\
Pin down the exact curve by plugging in our starting point $(0, 1)$:
$$1 = Ce^0 \implies C = 1$$

\textbf{Final Unique Solution:}
$y(x) = e^x$
\end{frame}

\begin{frame}{Solving the BVP: The Uniqueness Trap}
\textbf{Problem:} $y'' + y = 0,$ bounded at $x=0$ and $x=\pi.$

\textbf{Step 1: Find the General Solution} \\
The characteristic equation gives roots $\pm i$, yielding sines and cosines:
$$y(x) = c_1 \cos(x) + c_2 \sin(x)$$

\textbf{Step 2: Apply the First Boundary ($y(0) = 0$)}
$$0 = c_1 \cos(0) + c_2 \sin(0) \implies c_1 = 0$$
Our equation shrinks to $$y(x) = c_2 \sin(x)$$.

\textbf{Step 3: The Second Boundary Dictates Solvability}
\begin{itemize}
    \item \textbf{If} $y(\pi) = 0$: $0 = c_2 \sin(\pi) \implies 0 = 0$. $c_2$ can be anything! \\ \textit{Result: Infinitely many solutions (e.g., any standing wave on a string).}
    \item \textbf{If} $y(\pi) = 1$: $1 = c_2 \sin(\pi) \implies 1 = 0$. Contradiction! \\ \textit{Result: Absolutely zero solutions.}
\end{itemize}
\end{frame}


\begin{frame}{Geometric Interpretation}
Differential equations define a \textbf{vector field}.

Example:

\[
\frac{dx}{dt} = f(x)
\]

Each point has a vector that determines the direction of motion.

Solutions are called \textbf{integral curves}.
\end{frame}


%\begin{frame}{Physical Motivation: Newton's Second Law}
%Newton's law:
%
%\[
%F = m\frac{d^2x}{dt^2}
%\]
%
%If the force depends on position:
%
%\[
%m\frac{d^2x}{dt^2} = F(x)
%\]
%
%This is a \textbf{second-order differential equation} describing motion.
%\end{frame}
%
%\begin{frame}{Example: Exponential Growth}
%Population growth model:
%
%\[
%\frac{dN}{dt} = rN
%\]
%
%Solution:
%
%\[
%N(t) = N_0 e^{rt}
%\]
%
%Interpretation: the growth rate is proportional to the current population.
%\end{frame}
%
%\begin{frame}{Example: RC Circuit}
%Kirchhoff's law gives
%
%\[
%RC\frac{dV}{dt} + V = V_{in}(t)
%\]
%
%This is a \textbf{first-order linear differential equation}.
%
%Such equations appear in many physical systems.
%\end{frame}

\begin{frame}{Direction Fields}
For a first-order equation

\[
\frac{dy}{dx} = f(x,y)
\]

Each point $(x,y)$ has a slope:

\[
slope = f(x,y)
\]

Plotting these slopes produces a \textbf{direction field}.

Solutions follow these directions.
\end{frame}

\begin{frame}{Phase Portraits}
For systems:

\[
\dot{x} = f(x,y), \quad \dot{y} = g(x,y)
\]

We can study trajectories in the $(x,y)$ plane.

This geometric picture is called a \textbf{phase portrait}.

It reveals stability, oscillations, and equilibrium points.
\end{frame}


%%%%%%%%%%%%%%%%%%%%%%%%%%%%%%%%%%%%%%%%%%%%%%%%%%%%%%%%%%%%%%%%%%%%%%%%%%%%


\begin{frame}{What Does $\frac{dy}{dx}$ Actually Mean?}

\[
\frac{dy}{dx} = f(x,y)
\]

can be written as a differential form:

\[
dy - f(x,y)dx = 0
\]

Solutions are curves along which this 1-form vanishes.

This connects differential equations with

\begin{itemize}
\item differential forms
\item exterior derivatives
\item integrability conditions
\end{itemize}
\end{frame}

\begin{frame}{Vector Fields}
A first-order equation

\[
\frac{dx}{dt} = f(x)
\]

defines a \textbf{vector field}.

Each point $x$ has an associated velocity vector $f(x)$.

Solutions are curves tangent to the vector field.

These curves are called \textbf{integral curves}.
\end{frame}

\begin{frame}{Flows on Manifolds}
Given a vector field $X$ on a manifold $M$:

\[
\dot{x} = X(x)
\]

its solutions define a \textbf{flow}

\[
\varphi_t : M \rightarrow M
\]

which moves points along integral curves.

Thus differential equations generate dynamical systems on manifolds.
\end{frame}

\begin{frame}{Tangent Vectors}
On a manifold the derivative defines a \textbf{tangent vector}:

\[
\frac{d}{dt}x(t) \in T_xM
\]

where $T_xM$ is the tangent space.

Differential equations describe curves whose velocities lie in these tangent spaces.

Hence they are naturally formulated using:

\begin{itemize}
\item tangent bundles
\item vector fields
\item smooth manifolds
\end{itemize}
\end{frame}

\begin{frame}{Connections Across Mathematics}
The simple symbol

\[
\frac{dy}{dx}
\]

connects many areas of mathematics:

\begin{itemize}
\item calculus
\item differential geometry
\item dynamical systems
\item differential forms
\item topology of flows
\item numerical analysis
\item physics (laws of motion)
\end{itemize}

Differential equations sit at the intersection of these fields.
\end{frame}


\begin{frame}{One Equation in Many Mathematical Languages}

Consider the differential equation

\[
\frac{dy}{dx} = f(x,y)
\]

The same object can be interpreted in several ways:

\begin{itemize}
\item \textbf{Slope field}: each point $(x,y)$ has slope $f(x,y)$
\item \textbf{Vector field}:
\[
V(x,y) = (1,f(x,y))
\]
\item \textbf{Differential form}:
\[
dy - f(x,y)\,dx = 0
\]
\item \textbf{Integral curves}: solutions are curves tangent to $V$
\item \textbf{Flow of a vector field}: trajectories generated by motion in the field
\item \textbf{Dynamical system}: evolution rule for a state
\end{itemize}

A differential equation is therefore a \textbf{geometric object}.
\end{frame}



\begin{frame}{Differential Equations as Geometry}

A useful hierarchy of ideas:

\[
\text{functions}
\]

\[
\downarrow
\]

\[
\text{derivatives}
\]

\[
\downarrow
\]

\[
\text{vector fields}
\]

\[
\downarrow
\]

\[
\text{flows on spaces}
\]

\[
\downarrow
\]

\[
\text{dynamical systems}
\]

Thus differential equations describe \textbf{motion in geometric spaces}.

These spaces may be

\begin{itemize}
\item $\mathbb{R}^n$
\item manifolds
\item configuration spaces in physics
\item phase spaces
\end{itemize}

This viewpoint connects differential equations with
\textbf{geometry, topology, and physics}.
\end{frame}



\begin{frame}{From Curves to Manifolds}

A curve

\[
x(t)
\]

has velocity

\[
\dot{x}(t)
\]

which is a vector in the \textbf{tangent space}

\[
\dot{x}(t) \in T_{x(t)}M
\]

Thus a differential equation

\[
\dot{x} = X(x)
\]

specifies a \textbf{vector field}

\[
X : M \rightarrow TM
\]

and solutions are curves whose velocity matches the field.

This is the geometric formulation of differential equations on manifolds.
\end{frame}



\begin{frame}{Unifying Picture}

\[
\frac{dy}{dx}
\]

can simultaneously be seen as

\begin{itemize}
\item a \textbf{ratio of infinitesimal changes}
\item a \textbf{limit}
\item a \textbf{tangent slope}
\item a \textbf{linear approximation}
\item a \textbf{vector field component}
\item a \textbf{differential form relation}
\item a \textbf{generator of a flow}
\end{itemize}

This simple symbol sits at the intersection of

\begin{center}
calculus \quad geometry \quad topology \quad physics
\end{center}

Differential equations are one of the main bridges between these fields.
\end{frame}


\end{document}

